% GENERATED FILE. Edit canonical lessons, then run npm run generate:books.
% canonical-source-hash: fnv1a64:2495a4e8c6f2bbfa
% canonical-lessons: IT-C349-effetto, IT-C349-provocare, IT-C349-causare, IT-C349-fattore, IT-C349-ragione

\chapter{Cause ed effetti}
\label{ch:it-cause-ed-effetti}
\hlchaptermodality{349}

\begin{chapteropening}
\textbf{By the end of this chapter:} \emph{I can talk about causes: say what caused or provoked something, name a factor, an effect and the reason behind it.}

Complete the last lesson of chapter 349: I can talk about causes: say what caused or provoked something, name a factor, an effect and the reason behind it.
\end{chapteropening}

\section[l'effetto]{l'effetto — the effect}

\label{lesson:IT-C349-effetto}



\begin{quote}
Before the new one: say the Italian for the criticism, then the Italian for to reiterate.
\end{quote}

\subsection*{You'll want to know: l'effetto}
\textbf{l'effetto} — \textquotedblleft{}the effect\textquotedblright{}.

\begin{quote}
La medicina ha fatto effetto dopo dieci minuti.
\end{quote}

\begin{quote}
\emph{The medicine took effect after ten minutes.}
\end{quote}

\emph{Fare effetto} — to take effect. \emph{In effetti} — in fact, indeed.

\begin{cousinweb}
\emph{effetto} comes from Latin \emph{effectus} \textquotedblleft{}an effect, a result, a purpose\textquotedblright{} (EtymDB), from \emph{efficere} \textquotedblleft{}to bring about\textquotedblright{}, built from \emph{ex} \textquotedblleft{}out\textquotedblright{} and \emph{facere} \textquotedblleft{}to do, to make\textquotedblright{}.
\end{cousinweb}

\begin{grammarlens}[title={What you've built}]
One more word for telling, explaining and describing at B1.
\end{grammarlens}

\subsection*{Your turn}
\begin{itemize}
\raggedright
  \item \emph{Say it:} \emph{l'effetto}
  \item \emph{Say it:} \emph{l'effetto}, once more
  \item \emph{Say it:} a sentence of your own with \emph{l'effetto}
  \item \emph{Recall:} say the Italian for the criticism, then the Italian for to reiterate, then say \emph{l'effetto} again
\end{itemize}

\begin{tcolorbox}[breakable,colback=teal!4,colframe=teal!35!black,title={Before you move on}]
What does l'effetto mean? (\textquotedblleft{}The effect\textquotedblright{}.) Say it once more.
\end{tcolorbox}
\section[provocare]{provocare — to bring about, to cause; to provoke}

\label{lesson:IT-C349-provocare}



\begin{quote}
Before the new one: say the Italian for to reiterate, then the Italian for the effect.
\end{quote}

\subsection*{You'll want to know: provocare}
\textbf{provocare} — \textquotedblleft{}to bring about, to cause; to provoke\textquotedblright{}.

\begin{quote}
La pioggia ha provocato molti incidenti.
\end{quote}

\begin{quote}
\emph{The rain caused many accidents.}
\end{quote}

With bad events it means to cause; with a person, to provoke: \emph{non provocarmi!}

\begin{cousinweb}
\emph{provocare} is a learned borrowing from Latin \emph{prōvocāre} \textquotedblleft{}to call forth, to challenge\textquotedblright{} (Lewis \& Short), built from \emph{prō} \textquotedblleft{}forth\textquotedblright{} and \emph{vocāre} \textquotedblleft{}to call\textquotedblright{}.
\end{cousinweb}

\begin{grammarlens}[title={What you've built}]
One more word for telling, explaining and describing at B1.
\end{grammarlens}

\subsection*{Your turn}
\begin{itemize}
\raggedright
  \item \emph{Say it:} \emph{provocare}
  \item \emph{Say it:} \emph{provocare}, once more
  \item \emph{Say it:} a sentence of your own with \emph{provocare}
  \item \emph{Recall:} say the Italian for to reiterate, then the Italian for the effect, then say \emph{provocare} again
\end{itemize}

\begin{tcolorbox}[breakable,colback=teal!4,colframe=teal!35!black,title={Before you move on}]
What does provocare mean? (\textquotedblleft{}To bring about, to cause; to provoke\textquotedblright{}.) Say it once more.
\end{tcolorbox}
\section[causare]{causare — to cause}

\label{lesson:IT-C349-causare}



\begin{quote}
Before the new one: say the Italian for the effect, then the Italian for to bring about.
\end{quote}

\subsection*{You'll want to know: causare}
\textbf{causare} — \textquotedblleft{}to cause\textquotedblright{}.

\begin{quote}
Il fumo può causare molte malattie.
\end{quote}

\begin{quote}
\emph{Smoking can cause many illnesses.}
\end{quote}

Neutral and frequent; the phrase \emph{a causa di} means because of.

\begin{cousinweb}
\emph{causare} is built on \emph{causa} \textquotedblleft{}cause\textquotedblright{}, a learned borrowing from Latin \emph{causa} \textquotedblleft{}cause, reason, lawsuit\textquotedblright{} (EtymDB). The everyday word \emph{cosa} \textquotedblleft{}thing\textquotedblright{} is the same Latin word, inherited and worn down.
\end{cousinweb}

\begin{grammarlens}[title={What you've built}]
One more word for telling, explaining and describing at B1.
\end{grammarlens}

\subsection*{Your turn}
\begin{itemize}
\raggedright
  \item \emph{Say it:} \emph{causare}
  \item \emph{Say it:} \emph{causare}, once more
  \item \emph{Say it:} a sentence of your own with \emph{causare}
  \item \emph{Recall:} say the Italian for the effect, then the Italian for to bring about, then say \emph{causare} again
\end{itemize}

\begin{tcolorbox}[breakable,colback=teal!4,colframe=teal!35!black,title={Before you move on}]
What does causare mean? (\textquotedblleft{}To cause\textquotedblright{}.) Say it once more.
\end{tcolorbox}
\section[il fattore]{il fattore — the factor}

\label{lesson:IT-C349-fattore}



\begin{quote}
Before the new one: say the Italian for to bring about, then the Italian for to cause.
\end{quote}

\subsection*{You'll want to know: il fattore}
\textbf{il fattore} — \textquotedblleft{}the factor\textquotedblright{}.

\begin{quote}
Il prezzo è un fattore importante nella scelta.
\end{quote}

\begin{quote}
\emph{The price is an important factor in the choice.}
\end{quote}

\emph{Il fattore tempo} — the time factor. In the country, \emph{il fattore} was also the farm manager.

\begin{cousinweb}
\emph{fattore} comes from Latin \emph{factor} \textquotedblleft{}maker, doer, performer\textquotedblright{} (EtymDB), built on \emph{facere} \textquotedblleft{}to make\textquotedblright{}: one of the things that make a result.
\end{cousinweb}

\begin{grammarlens}[title={What you've built}]
One more word for telling, explaining and describing at B1.
\end{grammarlens}

\subsection*{Your turn}
\begin{itemize}
\raggedright
  \item \emph{Say it:} \emph{il fattore}
  \item \emph{Say it:} \emph{il fattore}, once more
  \item \emph{Say it:} a sentence of your own with \emph{il fattore}
  \item \emph{Recall:} say the Italian for to bring about, then the Italian for to cause, then say \emph{il fattore} again
\end{itemize}

\begin{tcolorbox}[breakable,colback=teal!4,colframe=teal!35!black,title={Before you move on}]
What does il fattore mean? (\textquotedblleft{}The factor\textquotedblright{}.) Say it once more.
\end{tcolorbox}
\section[la ragione]{la ragione — the reason; reason (the faculty)}

\label{lesson:IT-C349-ragione}



\begin{quote}
Before the new one: say the Italian for to cause, then the Italian for the factor.
\end{quote}

\subsection*{You'll want to know: la ragione}
\textbf{la ragione} — \textquotedblleft{}the reason; reason (the faculty)\textquotedblright{}.

\begin{quote}
Per quale ragione sei arrivato tardi?
\end{quote}

\begin{quote}
\emph{For what reason did you arrive late?}
\end{quote}

\emph{Avere ragione} — to be right; \emph{dare ragione a qualcuno} — to say someone is right.

\begin{cousinweb}
\emph{ragione} is inherited from Latin \emph{ratiō, ratiōnis} \textquotedblleft{}a reckoning, an account, reason\textquotedblright{} (Lewis \& Short), built on \emph{rērī} \textquotedblleft{}to think, to reckon\textquotedblright{}.
\end{cousinweb}

\begin{grammarlens}[title={What you've built}]
One more word for telling, explaining and describing at B1.
\end{grammarlens}

\subsection*{Your turn}
\begin{itemize}
\raggedright
  \item \emph{Say it:} \emph{la ragione}
  \item \emph{Say it:} \emph{la ragione}, once more
  \item \emph{Say it:} a sentence of your own with \emph{la ragione}
  \item \emph{Recall:} say the Italian for to cause, then the Italian for the factor, then say \emph{la ragione} again
\end{itemize}

\begin{tcolorbox}[breakable,colback=teal!4,colframe=teal!35!black,title={Before you move on}]
What does la ragione mean? (\textquotedblleft{}The reason; reason (the faculty)\textquotedblright{}.) Say it once more.
\end{tcolorbox}
